\section*{Capítulo \ref{ch:g8:fronts-and-weather-maps} --- Massas de ar, frentes e mapas do tempo}

\begin{solution}{exo:g8:fronts-and-weather-maps:1}
(a) Tropical marítima: quente e úmida. (b) Polar continental: fria e
seca. (c) Tropical continental: muito quente e seca. (d) Polar marítima:
fria e úmida.
\end{solution}

\begin{solution}{exo:g8:fronts-and-weather-maps:2}
Linha azul com triângulos: uma \omterm{def:g8:fronts-and-weather-maps:front}{frente fria}. Linha vermelha com
semicírculos: uma \omterm{def:g8:fronts-and-weather-maps:front}{frente quente}. Linha que liga lugares com a mesma
pressão: uma \omterm{def:g8:fronts-and-weather-maps:isobar}{isóbara}. H: o centro de um \omterm{def:g8:fronts-and-weather-maps:depression}{anticiclone} (alta pressão). L: o
centro de uma \omterm{def:g8:fronts-and-weather-maps:depression}{depressão} (baixa pressão).
\end{solution}

\begin{solution}{exo:g8:fronts-and-weather-maps:3}
(a) Verdadeiro. (b) Falso: elas chegam antes dela, bem à frente, lá no
alto do céu. (c) Falso: \omterm{def:g8:fronts-and-weather-maps:isobar}{isóbaras} próximas significam vento forte.
(d) Verdadeiro.
\end{solution}

\begin{solution}{exo:g8:fronts-and-weather-maps:4}
(a) \qty{980}{hPa}. (b) \qty{1034}{hPa}. (c) \qty{4}{hPa}. (d) $1034 -
980 = \qty{54}{hPa}$.
\end{solution}

\begin{solution}{exo:g8:fronts-and-weather-maps:5}
Finos riscos de \omterm{def:g4:clouds-rain-snow:cloud-families}{cirros}; um véu cinzento de altostratos; as \omterm{def:g4:clouds-rain-snow:cloud}{nuvens} ficam
mais baixas e mais grossas; \omterm{def:g4:clouds-rain-snow:precipitation}{chuva} contínua; ar mais quente e mais úmido.
\end{solution}

\begin{solution}{exo:g8:fronts-and-weather-maps:6}
Numa \omterm{def:g8:fronts-and-weather-maps:front}{frente fria}, o ar frio e denso avança por baixo do ar quente como uma
\emph{cunha}: o ar quente é levantado numa encosta \emph{íngreme}, sobe
depressa e forma \omterm{def:g4:clouds-rain-snow:cloud}{nuvens} altas. Numa \omterm{def:g8:fronts-and-weather-maps:front}{frente quente}, o ar quente sobe
devagar por cima do ar frio numa \emph{rampa suave}: ele sobe pouco a
pouco ao longo de uma grande distância, e as suas \omterm{def:g4:clouds-rain-snow:cloud}{nuvens} se espalham em
camadas.
\end{solution}

\begin{solution}{exo:g8:fronts-and-weather-maps:7}
(a) Perto de L: as \omterm{def:g8:fronts-and-weather-maps:isobar}{isóbaras} estão mais próximas ali, e as setas são mais
longas. (b) Do leste (as setas apontam para o oeste, um pouco para o
sudoeste). (c) Do leste (as setas apontam para o oeste). (d) No sentido
anti-horário.
\end{solution}

\begin{solution}{exo:g8:fronts-and-weather-maps:8}
(a) $1016 - 1003 = \qty{13}{hPa}$. (b) Uma \omterm{def:g8:fronts-and-weather-maps:depression}{depressão}: a pressão continua
caindo. (c) \omterm{def:g4:clouds-rain-snow:cloud}{Nuvens}, vento e \omterm{def:g4:clouds-rain-snow:precipitation}{chuva}, provavelmente com uma \omterm{def:g8:fronts-and-weather-maps:front}{frente}.
\end{solution}

\begin{solution}{exo:g8:fronts-and-weather-maps:9}
(a) cT, tropical continental. (b) Muito quente e seca. (c) Sobre o \omterm{def:g2:water-on-earth:ocean}{mar},
ela recolhe vapor de água: fica mais úmida, e o ar perto da água é um
pouco resfriado pelo \omterm{def:g2:water-on-earth:ocean}{mar}.
\end{solution}

\begin{solution}{exo:g8:fronts-and-weather-maps:10}
(a) $4 - 1 = 3$ dias. (b) $3 \div 3 = 1$ dia a cada dez anos. (c) Ela
precisa de mais $5 - 1 = 4$ dias: cerca de 40 anos.
\end{solution}

\begin{solution}{exo:g8:fronts-and-weather-maps:11}
(a) Cerca de \qty{1350}{km} (quatro barras e meia de \omterm{def:g2:drawing-to-map:scale}{escala}). (b) \qty{4}{hPa}
em \qty{80}{km}: $4 \div 80 \times 100 = \qty{5}{hPa}$ a cada
\qty{100}{km}. (c) $4 \div 140 \times 100 \approx \qty{2.9}{hPa}$ a cada
\qty{100}{km}. (d) Perto de L, onde a pressão muda mais depressa; sim, as
setas são mais longas perto de L.
\end{solution}

\begin{solution}{exo:g8:fronts-and-weather-maps:12}
(a) A \omterm{def:g8:fronts-and-weather-maps:front}{frente fria}: $300 \div 50 = 6$ horas. A \omterm{def:g8:fronts-and-weather-maps:front}{frente quente} se desloca a
\qty{25}{km} por hora: $200 \div 25 = 8$ horas. A \omterm{def:g8:fronts-and-weather-maps:front}{frente fria} chega
2 horas antes. (b) Primeira cidade: \omterm{def:g4:clouds-rain-snow:cloud}{nuvens} altíssimas, \omterm{def:g4:clouds-rain-snow:precipitation}{chuva} forte,
rajadas, e depois ar mais fresco e céu abrindo. Segunda cidade: \omterm{def:g4:clouds-rain-snow:cloud}{nuvens} em
camadas que vão baixando e uma \omterm{def:g4:clouds-rain-snow:precipitation}{chuva} longa e contínua, e depois ar mais
quente e mais úmido. (c) A segunda cidade, alcançada pelo ar quente.
\end{solution}

\begin{solution}{pb:g8:fronts-and-weather-maps:1}
\textbf{1.} No ar quente, entre as duas \omterm{def:g8:fronts-and-weather-maps:front}{frentes}. A \omterm{def:g8:fronts-and-weather-maps:front}{frente quente} já
passou.

\textbf{2.} Do sudoeste.

\textbf{3.} Para o leste: os triângulos de uma \omterm{def:g8:fronts-and-weather-maps:front}{frente fria} ficam do lado
para onde ela se desloca, e eles apontam para o leste, em direção à cidade.

\textbf{4.} \omterm{def:g4:clouds-rain-snow:cloud}{Nuvens} altíssimas, \omterm{def:g4:clouds-rain-snow:precipitation}{chuva} forte, talvez trovões e \omterm{def:g4:clouds-rain-snow:precipitation}{granizo},
rajadas de vento; depois, ar mais fresco, vento de oeste ou de noroeste e
céu abrindo.

\textbf{5.} A \qty{180}{km} às 06:00 e a \qty{300}{km} às 09:00.

\textbf{6.} $300 - 180 = \qty{120}{km}$.

\textbf{7.} $120 \div 3 = \qty{40}{km}$ por hora.

\textbf{8.} A cidade fica a \qty{480}{km}: $480 - 300 = \qty{180}{km}$.

\textbf{9.} $180 \div 40 = 4.5$ horas, ou seja, 4 horas e 30 minutos.

\textbf{10.} 09:00 + 4~h~30~min = \textbf{13:30}.

\textbf{11.} Não: as pancadas devem terminar por volta das 14:30, antes do
começo do jogo.

\textbf{12.} $180 \div 30 = 6$ horas: às 15:00, bem no começo do jogo.
Previsão: ``A \omterm{def:g8:fronts-and-weather-maps:front}{frente fria} deve chegar à cidade por volta das
\textbf{13:30}, com uma hora de \omterm{def:g4:clouds-rain-snow:precipitation}{chuva} forte e rajadas, e depois o céu vai
abrir: o jogo provavelmente poderá ser disputado. Se a \omterm{def:g8:fronts-and-weather-maps:front}{frente} desacelerar,
ela poderá chegar às 15:00 e atrasar o jogo.''
\end{solution}
